% =====================================================================
% IV. CONCLUSION
% LLM declaration: text level 3; see Appendix "Use of AI/LLM tools".
% =====================================================================

\section{Conclusion}\label{sec:conclusion}

For OLS the test error of Runge's function falls with the polynomial degree until the variance takes over, and where that happens depends on the data: the best degree rose from 6 to 20 when $n$ went from 30 to 1000, and fell when the noise increased. The largest errors came from a few data sets where the model had to extrapolate or bridge a gap near the edges, so a single split, or the mean over many, can give a misleading picture. Ridge shrinks the directions of the design matrix with small singular values, which here contain mostly noise. With $n = 100$ it did not improve on OLS at the best fixed degree, but it removed the blow-up at high degree, and with $n = 30$ it gave the lowest error. The bootstrap split the error into bias and variance, with the variance overtaking the measured bias from degree 11 for $n = 100$, but it overestimated the error at high degree because each bootstrap sample contains only about 63\,\% distinct points; cross-validation followed the true error closely. In the final comparison by cross-validation, Ridge with a small penalty was best (0.0149 against 0.0158 for OLS), mainly because it keeps the variance down at high degree, but the differences between the methods were small compared with the differences between data sets. For the optimisation, the largest usable learning rate of plain gradient descent is set by the steepest curvature and the number of iterations by the flattest. Momentum gave the largest speed-up (601 against 11\,834 iterations), and Adam reached the criterion for every learning rate we tested, while RMSProp and Adam with a constant learning rate did not settle exactly at the minimum. Gradient descent with subgradients approached the Lasso solution when the step was small enough, but never gave exact zeros. With only 80 training points, stochastic gradient descent reached a moderate accuracy about as fast as full-batch gradient descent but did not settle as low, so it gave no clear advantage.

Several results are specific to our setup: one function in one dimension, uniformly drawn $x$ values, and one noise model. Many of our references, such as the error on new data and the bias measured against $f$, are possible only because $f$ is known. The resampling studies used 20 data sets, and the gradient-descent experiments one data set at degree 6, so the exact numbers there carry noticeable uncertainty. The median, which we used throughout, hides the extreme cases that a single real data set may well be. At high degree the bootstrap estimate of the bias is contaminated by a few extreme fits. Finally, scikit-learn's Lasso solver did not converge for small penalties at high degree, so our conclusions about Lasso with small penalties are uncertain; with the penalties where it converged, Lasso was clearly worse than Ridge.

Natural next steps would be to sample $x$ at Chebyshev nodes, which are denser near the edges, or to use an orthogonal polynomial basis, to reduce both the oscillations and the collinearity of the columns; to use a Lasso solver that converges for small penalties, to settle the comparison with Ridge; and to test decaying learning rates for Adam and stochastic gradient descent on a problem with many more data points, where stochastic gradient descent is expected to pay off.
